\section{Conclusion and Future Work}
\label{sec:conclusion}

In this project, we built and evaluated a leakage-free machine learning pipeline to predict First Payment Default across $233{,}154$ LTFS vehicle loans, yielding four main conclusions:

\begin{enumerate}[itemsep=1pt,topsep=2pt]
  \item \textbf{Gradient Boosting Leads Credit Risk Separation:} Tuned \texttt{LightGBM} ($\text{OOT AUC} = 0.6621$) and \texttt{XGBoost} ($0.6620$) achieve top performance and clearly outperform \texttt{RandomForest} ($0.6554$), \texttt{LogisticRegression} ($0.6525$), the controlled \texttt{MLP} baseline ($0.6525$), and \texttt{DecisionTree} ($0.6270$). \texttt{LightGBM} provides the best balance of training speed ($10.9\text{ s}$), Development calibration ($\text{ECE} = 0.0017$), and exact TreeSHAP explainability.
  \item \textbf{Domain Engineering and Channel Encoding Matter:} Separating non-scored credit bureau codes ($<300$) from numeric scores ($300\text{--}890$) and adding leverage/distress ratios improves OOT generalization across all models ($+0.0021$ to $+0.0026$ AUC), while leak-free Out-of-Fold Empirical-Bayes target encoding on dealer, branch, and postal channels contributes $+0.0115$ to $+0.0180$ AUC gain on future cohorts.
  \item \textbf{Time-Based Validation Prevents Optimism Bias:} Standard random 80/20 splitting overstates OOT ROC-AUC by $+0.0091$ to $+0.0106$ and conceals a $30$-fold calibration degradation ($\text{ECE } 0.0016 \to 0.0491$) triggered when late-October default prevalence jumps from $20.63\%$ to $25.59\%$.
  \item \textbf{Threshold Optimization Beats Synthetic Resampling:} Fold-isolated \texttt{SMOTE} degrades \texttt{LightGBM} OOT ROC-AUC down to $0.6252$, and heavy balanced class weighting distorts calibration ($\text{ECE} = 0.2111$). Training on natural empirical priors coupled with a Bayesian $5{:}1$ cost threshold ($\tau = 0.167$) preserves ranking while nearly halving principal-weighted portfolio loss compared to the naive $0.50$ cutoff ($7.13\%$ vs.\ $13.12\%$).
\end{enumerate}

\subsection{Limitations and Future Work}
Four methodological and data boundaries should be considered when deploying these findings. First, the labeled corpus spans three months (August--October 2018) and omits borrower income, debt-to-income ratios, loan tenure, and vehicle depreciation, creating an observable information ceiling near $\text{OOT AUC} \approx 0.662$ ($0.645$ on thin-file loans); incorporating alternative cashflow or utility signals could narrow this gap. Second, the neural \texttt{MLP} was evaluated as a controlled baseline architecture without exhaustive Optuna hyperparameter tuning or learned entity embeddings, so deeper tabular architectures could be revisited under equalized tuning budgets. Third, nearly $29\%$ of OOT SHAP attribution stems from dealer and postal proxies, motivating fair-lending constraints and rolling target-encoding refreshes as channels evolve. Finally, production deployment should incorporate online log-odds intercept recalibration for seasonal prior shifts and reject-inference auditing across multi-quarter horizons.
